%======================================================================
% Chapter 11: Mathematics for Deep Learning
%======================================================================
\chapter{Mathematics for Deep Learning}\label{ch:dl}

\begin{motivation}
A deep neural network is a long composition of simple functions: matrix
multiplications followed by simple non-linear functions. It turns its input
into probabilities, is scored by a loss derived from probability theory, and is
trained by gradient descent, with gradients computed by the chain rule. Every
one of these ingredients has appeared in this book:
\begin{itemize}
  \item layers are matrix--vector products (\cref{ch:linalg1,ch:linalg2});
  \item activation functions, exponentials and logarithms come from
    \cref{ch:functions};
  \item gradients and the chain rule come from \cref{ch:calculus,ch:multivariable};
  \item training is optimisation (\cref{ch:optimization});
  \item outputs are probability distributions, losses are negative
    log-likelihoods, and generalisation is a statistical question
    (\cref{ch:probability,ch:statistics}).
\end{itemize}
This final technical chapter puts them together. We study softmax outputs, the
information theory behind the cross-entropy loss, the matrix calculus that
backpropagation needs, backpropagation itself, the problem of vanishing
gradients, automatic differentiation and attention, and we end with a complete
neural network trained from scratch.
\end{motivation}

\begin{objectives}
\begin{itemize}
  \item describe a feed-forward network as a composition of layers and count its
    parameters;
  \item compute softmax probabilities, and do so in a numerically stable way;
  \item compute entropy, cross-entropy and KL divergence, and explain why
    cross-entropy is the loss for classification;
  \item use the basic identities of matrix calculus and check them with shapes;
  \item derive the gradient $\vq - \vy$ of the softmax cross-entropy loss;
  \item carry out backpropagation by hand on a small network;
  \item explain vanishing and exploding gradients and their remedies;
  \item explain how automatic differentiation works;
  \item compute scaled dot-product attention;
  \item assemble the complete training loop of a neural network.
\end{itemize}
\end{objectives}

%----------------------------------------------------------------------
\section{Neural networks: compositions of simple functions}\label{sec:neural-networks}
%----------------------------------------------------------------------

A linear model $\hat y = \vw^\top\vx + b$ can only separate classes by a straight
line (or a plane, in higher dimensions). Neural networks get their flexibility
by stacking many simple non-linear steps.

\begin{definition}[Neuron and layer]\label{def:layer}
A \term{neuron} computes a weighted sum of its inputs plus a bias, the
\term{pre-activation} $z = \vw^\top\vx + b$, and outputs $a = \varphi(z)$, where
$\varphi$ is an \term{activation function} such as ReLU, sigmoid or $\tanh$
(\cref{sec:activation-gallery}). A \emph{layer} of $m$ neurons with $n$ inputs
computes
\[
  \vz = \mW\vx + \vb, \qquad \va = \varphi(\vz),
\]
where $\mW$ is an $m\times n$ matrix (one row of weights per neuron),
$\vb \in \R^m$, and $\varphi$ is applied to each component.
\end{definition}

\begin{definition}[Feed-forward network]\label{def:mlp}
A \term{feed-forward neural network} (or \emph{multilayer perceptron}) with $L$
layers computes, starting from $\va^{(0)} = \vx$,
\[
  \keyeq{\vz^{(l)} = \mW^{(l)}\va^{(l-1)} + \vb^{(l)}, \qquad
  \va^{(l)} = \varphi\bigl(\vz^{(l)}\bigr), \qquad l = 1, \dots, L}
\]
The superscript $(l)$ numbers the layers. The layers between input and output
are called \term[hidden layer]{hidden layers}. The output layer usually has no
activation, or a special one (\cref{sec:softmax}), and its outputs
$\vz^{(L)}$ are called \term{logits}.
\end{definition}

\begin{figure}[htbp]
\centering
\begin{tikzpicture}[>={Stealth[length=2mm]},
  neuron/.style={circle,draw,thick,minimum size=8mm,inner sep=0pt},
  innode/.style={neuron,fill=MLdef!12},
  hidnode/.style={neuron,fill=MLexample!14},
  outnode/.style={neuron,fill=MLintuition!22}]
  \foreach \i in {1,2,3}{\node[innode] (x\i) at (0,2.4-1.2*\i) {$x_\i$};}
  \foreach \i in {1,...,4}{\node[hidnode] (h\i) at (4,3-1.2*\i) {};}
  \foreach \i in {1,2}{\node[outnode] (o\i) at (8,1.8-1.2*\i) {};}
  \foreach \i in {1,2,3}{\foreach \j in {1,...,4}{\draw[->,gray!70] (x\i) -- (h\j);}}
  \foreach \i in {1,...,4}{\foreach \j in {1,2}{\draw[->,gray!70] (h\i) -- (o\j);}}
  \node[font=\small,align=center,anchor=north] at (0,-2.5) {input\\$\vx \in \R^3$};
  \node[font=\small,align=center,anchor=north] at (4,-2.5) {hidden layer\\$\va^{(1)} = \varphi(\mW^{(1)}\vx + \vb^{(1)})$};
  \node[font=\small,align=center,anchor=north] at (8.2,-2.5) {output\\$\vz^{(2)} = \mW^{(2)}\va^{(1)} + \vb^{(2)}$};
  \node[font=\small,MLdef] at (2,2.6) {$\mW^{(1)}$: $4\times3$};
  \node[font=\small,MLdef] at (6,2.6) {$\mW^{(2)}$: $2\times4$};
\end{tikzpicture}
\caption{A network with $3$ inputs, one hidden layer of $4$ neurons and $2$
  outputs. Each arrow carries one weight; each neuron also has a bias.}
\label{fig:network}
\end{figure}

\begin{example}[A two-layer network]\label{ex:two-layer-forward}
Take the layer of \cref{sec:matrix-vector}: $\mW^{(1)} = \left(\begin{smallmatrix} 1 &
-1 & 2\\ 0 & -2 & 1\end{smallmatrix}\right)$, $\vb^{(1)} = (0.5, -1)$, ReLU
activation and input $\vx = (1, 2, 3)$, which gave $\va^{(1)} = (5.5, 0)$. Add an
output layer with $\mW^{(2)} = (1\ \ 3)$ and $b^{(2)} = -2$. The network's output
is $z^{(2)} = 1\cdot5.5 + 3\cdot0 - 2 = 3.5$. It has $2\cdot3 + 2 = 8$ parameters in
the first layer and $2 + 1 = 3$ in the second.
\end{example}

In general a layer with $n$ inputs and $m$ outputs has $mn + m$ parameters. A
classic network for recognising handwritten digits, with $784$ inputs (the
pixels of a $28\times28$ image), $128$ hidden units and $10$ outputs, has
$784\cdot128 + 128 + 128\cdot10 + 10 = 101\,770$ parameters; large language
models have hundreds of billions.

\begin{proposition}[Depth needs non-linearity]\label{prop:linear-collapse}
Without activation functions, a network collapses to a single linear layer:
\[
  \mW^{(2)}\bigl(\mW^{(1)}\vx + \vb^{(1)}\bigr) + \vb^{(2)}
  = \bigl(\mW^{(2)}\mW^{(1)}\bigr)\vx + \bigl(\mW^{(2)}\vb^{(1)} + \vb^{(2)}\bigr).
\]
\end{proposition}

This is \cref{prop:compose-linear} again, in matrix form: compositions of
linear functions are linear, so stacking linear layers adds nothing. The
activation functions are what make depth useful. A network of ReLU layers
computes a function that is linear on each of many small regions, with kinks
between them; with enough neurons, such functions can follow any shape.

\begin{remark}[Universal approximation]
A network with a single hidden layer can approximate any continuous function on
a bounded region as accurately as desired, provided it has enough hidden
neurons. This \emph{universal approximation theorem} guarantees that neural
networks are flexible enough; it says nothing about how many neurons are needed
or how to find the weights. In practice, deep networks represent complicated
functions far more economically than wide, shallow ones.
\end{remark}

\sectionexercises

\begin{exercise}
How many parameters does a network with $2$ inputs, one hidden layer of $3$
neurons and $1$ output have? And one with layers of sizes $10$, $20$, $20$, $5$?
\end{exercise}
\begin{solution}
$2\cdot3 + 3 + 3\cdot1 + 1 = 13$. For $10$--$20$--$20$--$5$:
$(200 + 20) + (400 + 20) + (100 + 5) = 745$.
\end{solution}

\begin{exercise}
Two linear layers have $\mW^{(1)} = \left(\begin{smallmatrix} 1 & 2\\ 0 &
1\end{smallmatrix}\right)$, $\vb^{(1)} = (1, 0)$, $\mW^{(2)} = (1\ \ {-1})$ and
$b^{(2)} = 2$, and no activation. Write the network as a single linear function
$w_1x_1 + w_2x_2 + c$.
\end{exercise}
\begin{solution}
$\mW^{(2)}\mW^{(1)} = (1\ \ 1)$ and $\mW^{(2)}\vb^{(1)} + b^{(2)} = 1 + 2 = 3$, so the
network computes $x_1 + x_2 + 3$.
\end{solution}

\begin{exercise}
In \cref{ex:two-layer-forward}, change the input to $\vx = (0, 1, 0)$ and compute
the output.
\end{exercise}
\begin{solution}
$\mW^{(1)}\vx + \vb^{(1)} = (-1 + 0.5,\ -2 - 1) = (-0.5, -3)$, so $\va^{(1)} = (0,
0)$ and the output is $z^{(2)} = 0 + 0 - 2 = -2$.
\end{solution}

%----------------------------------------------------------------------
\section{From scores to probabilities: the softmax function}\label{sec:softmax}
%----------------------------------------------------------------------

A classifier's last layer produces arbitrary real numbers, the logits. To turn
them into probabilities we need positive numbers that add up to $1$.

For \textbf{two classes}, one logit suffices: $q = \sigma(z) = \frac{1}{1 + e^{-z}}$
is the probability of class $1$, as in logistic regression
(\cref{sec:losses-from-likelihood}). The logit is then the log-odds,
$z = \ln\frac{q}{1 - q}$, the inverse of the sigmoid (\cref{sec:inverse-functions}).

For \textbf{$K$ classes}, the network outputs $K$ logits, and we exponentiate
them (making them positive, \cref{sec:exponential}) and normalise.

\begin{definition}[Softmax]\label{def:softmax}
The \term{softmax} of a vector $\vz \in \R^K$ is the vector $\vq \in \R^K$ with
\[
  \keyeq{q_k = \softmax(\vz)_k = \frac{e^{z_k}}{\sum_{j=1}^K e^{z_j}}}
\]
\end{definition}

\begin{example}[Computing a softmax]\label{ex:softmax}
For the logits $\vz = (2, 1, 0.1)$: $e^2 \approx 7.389$, $e^1 \approx 2.718$ and
$e^{0.1} \approx 1.105$, with sum $11.212$. So
$\softmax(\vz) \approx (0.659, 0.242, 0.099)$. The largest logit gets the largest
probability, but the others are not ignored.
\end{example}

\begin{theorem}[Properties of softmax]\label{thm:softmax-properties}
For every $\vz \in \R^K$:
\begin{enumerate}
  \item the outputs are positive and add up to $1$, so they form a categorical
    distribution (\cref{sec:distributions});
  \item larger logits get larger probabilities;
  \item \textbf{shift invariance:} $\softmax(\vz + c\vone) = \softmax(\vz)$ for any
    number $c$, where $\vone = (1, \dots, 1)$;
  \item for $K = 2$, $\softmax(z_1, z_2)_1 = \sigma(z_1 - z_2)$: softmax
    generalises the sigmoid.
\end{enumerate}
\end{theorem}

\begin{proof}
(1) and (2) hold because $e^x$ is positive and increasing. (3):
$e^{z_k + c} = e^ce^{z_k}$, and the factor $e^c$ cancels between numerator and
denominator. (4): dividing numerator and denominator by $e^{z_1}$,
\[
  \frac{e^{z_1}}{e^{z_1} + e^{z_2}} = \frac{1}{1 + e^{-(z_1 - z_2)}}
  = \sigma(z_1 - z_2). \qedhere
\]
\end{proof}

\begin{intuition}[A soft maximum]
If the logits are multiplied by a large number, the largest one dominates and
softmax approaches the one-hot vector of the largest logit, a ``hard'' maximum.
Softmax is a smooth, differentiable version of picking the maximum, which is
exactly what gradient-based training needs.
\end{intuition}

\subsection{Temperature}

Dividing the logits by a \term{temperature} $T > 0$ before the softmax,
$\vq = \softmax(\vz/T)$, controls how peaked the distribution is
(\cref{fig:temperature}). A low temperature makes the model more decisive; a high
temperature makes it more uniform. Language models use temperature when
generating text: low temperatures give predictable text, high temperatures give
varied (and eventually incoherent) text.

\begin{figure}[htbp]
\centering
\begin{tikzpicture}
\begin{axis}[mlplot,axis lines=left,ybar,bar width=11pt,width=10cm,height=5.4cm,
  ymin=0,ymax=1,xtick={1,2,3},xticklabels={class 1 ($z = 2$),class 2 ($z = 1$),class 3 ($z = 0.1$)},
  enlarge x limits=0.25,ylabel={probability},
  ylabel style={at={(ticklabel cs:0.5)},anchor=south,rotate=90},
  legend style={font=\small,at={(0.97,0.97)},anchor=north east,draw=gray!40},
  legend cell align=left]
  \addplot[fill=PlotBlue!70,draw=PlotBlue] coordinates {(1,0.864) (2,0.117) (3,0.019)};
  \addlegendentry{$T = 0.5$}
  \addplot[fill=PlotGreen!60,draw=PlotGreen] coordinates {(1,0.659) (2,0.242) (3,0.099)};
  \addlegendentry{$T = 1$}
  \addplot[fill=PlotOrange!60,draw=PlotOrange] coordinates {(1,0.502) (2,0.304) (3,0.194)};
  \addlegendentry{$T = 2$}
\end{axis}
\end{tikzpicture}
\caption{The softmax of the logits $(2, 1, 0.1)$ at three temperatures. Lower
  temperatures sharpen the distribution; higher ones flatten it.}
\label{fig:temperature}
\end{figure}

\subsection{Computing softmax safely}

A computer cannot evaluate $e^{1000}$: it overflows (double-precision numbers
stop at about $1.8\times10^{308} \approx e^{709.8}$). Shift invariance is the
cure: subtract the largest logit first. For example,
$\softmax(1000, 999) = \softmax(1, 0) \approx (0.731, 0.269)$. The same trick
evaluates the \term{log-sum-exp} function, which appears in the logarithm of
the softmax:
\[
  \ln\sum_{j=1}^K e^{z_j} = m + \ln\sum_{j=1}^K e^{z_j - m}, \qquad m = \max_j z_j,
  \qquad\text{and}\qquad
  \ln q_k = z_k - \ln\sum_{j=1}^K e^{z_j} .
\]

\sectionexercises

\begin{exercise}
Compute $\softmax(0, \ln2, \ln3)$ exactly.
\end{exercise}
\begin{solution}
The exponentials are $1$, $2$ and $3$, with sum $6$, so the result is
$\left(\tfrac16, \tfrac13, \tfrac12\right)$.
\end{solution}

\begin{exercise}
Compute $\softmax(1000, 1001, 1002)$ by hand, using shift invariance.
\end{exercise}
\begin{solution}
It equals $\softmax(0, 1, 2)$. The exponentials are $1$, $2.718$ and $7.389$, with
sum $11.107$, giving approximately $(0.090, 0.245, 0.665)$.
\end{solution}

\begin{exercise}
What does $\softmax(\vz/T)$ approach as $T \to \infty$? As $T \to 0$ (assuming the
largest logit is unique)?
\end{exercise}
\begin{solution}
As $T \to \infty$ all the scaled logits approach $0$, and the output approaches
the uniform distribution $\left(\tfrac1K, \dots, \tfrac1K\right)$. As $T \to 0$ the
differences between the scaled logits blow up, and the output approaches the
one-hot vector of the largest logit.
\end{solution}

\begin{exercise}
Show that $\max_jz_j \le \ln\sum_{j=1}^Ke^{z_j} \le \max_jz_j + \ln K$. (This is why
log-sum-exp is called a ``soft maximum''.)
\end{exercise}
\begin{solution}
Let $m = \max_jz_j$. The sum $\sum_je^{z_j}$ contains the term $e^m$ and consists of
$K$ terms that are each at most $e^m$, so $e^m \le \sum_je^{z_j} \le Ke^m$. Taking
logarithms gives $m \le \ln\sum_je^{z_j} \le m + \ln K$.
\end{solution}

%----------------------------------------------------------------------
\section{Information theory: entropy, cross-entropy and KL divergence}\label{sec:information-theory}
%----------------------------------------------------------------------

Why is $-\ln q$ the right loss for classification? \Cref{ch:statistics} gave one
answer, maximum likelihood. Information theory, founded by Claude Shannon in
1948 to study communication, gives a second one, and with it a vocabulary
(entropy, cross-entropy, KL divergence) used throughout machine learning.

\subsection{Surprise and entropy}

Learning that an unlikely event happened is more informative than learning that
a likely one did. The \term{information content} or \emph{surprise} of an event
of probability $p$ is
\[
  I = -\log_2 p \quad\text{bits}.
\]
A fair coin toss carries $-\log_2\frac12 = 1$ bit of information; a roll of a fair
die $\log_26 \approx 2.585$ bits; an event of probability $\frac{1}{1024}$, $10$
bits. A certain event carries none. For independent events probabilities
multiply, so surprises add (\cref{thm:log-laws}).

\begin{definition}[Entropy]\label{def:entropy}
The \term{entropy} of a discrete distribution $\vp = (p_1, \dots, p_K)$ is its
expected surprise,
\[
  \keyeq{H(\vp) = -\sum_{k=1}^K p_k\log p_k}
\]
with the convention $0\log0 = 0$. With $\log_2$ it is measured in \emph{bits};
with the natural logarithm, as is usual in machine learning, in \emph{nats}.
\end{definition}

\begin{example}[Entropies]\label{ex:entropies}
A fair coin has entropy $1$ bit. A coin with $P(\text{heads}) = 0.9$ has
$-(0.9\log_20.9 + 0.1\log_20.1) \approx 0.469$ bits: it is much more predictable. A
fair die has $\log_26 \approx 2.585$ bits. The distribution
$\left(\frac12, \frac14, \frac18, \frac18\right)$ has
$\frac12\cdot1 + \frac14\cdot2 + \frac18\cdot3 + \frac18\cdot3 = 1.75$ bits.
\end{example}

\begin{figure}[htbp]
\centering
\begin{tikzpicture}
\begin{axis}[mlplot,axis lines=left,width=8.5cm,height=5.4cm,xmin=0,xmax=1,ymin=0,ymax=1.12,clip=false,
  xlabel={$p = P(\text{heads})$},ylabel={$H$ (bits)},
  xlabel style={at={(ticklabel cs:0.5)},anchor=north},
  ylabel style={at={(ticklabel cs:0.5)},anchor=south,rotate=90}]
  \addplot[PlotBlue,very thick,domain=0.0005:0.9995,samples=200]
    {-x*ln(x)/ln(2) - (1-x)*ln(1-x)/ln(2)};
  \addplot[only marks,mark=*,mark size=1.8pt,PlotRed] coordinates {(0.5,1) (0.9,0.469)};
  \node[font=\small,anchor=south] at (axis cs:0.5,1) {fair coin};
  \node[font=\small,anchor=north east] at (axis cs:0.9,0.469) {$0.469$};
\end{axis}
\end{tikzpicture}
\caption{The entropy of a coin with $P(\text{heads}) = p$. It is largest, $1$ bit,
  for a fair coin and $0$ when the outcome is certain.}
\label{fig:binary-entropy}
\end{figure}

\begin{intuition}[Entropy as the length of a code]
Entropy is the average number of yes/no questions (bits) needed to communicate
an outcome, using the best possible code. For
$\left(\frac12, \frac14, \frac18, \frac18\right)$ the code
$\texttt{0}, \texttt{10}, \texttt{110}, \texttt{111}$ gives frequent outcomes short
codewords, with average length $\frac12\cdot1 + \frac14\cdot2 + \frac18\cdot3 +
\frac18\cdot3 = 1.75$ bits, exactly the entropy. Shannon proved that no code can
do better on average. In general, the best code spends $-\log_2p_k$ bits on an
outcome of probability $p_k$.
\end{intuition}

\subsection{Cross-entropy and KL divergence}

Suppose the outcomes follow $\vp$, but we use a code designed for a different
distribution $\vq$, spending $-\log q_k$ on outcome $k$.

\begin{definition}[Cross-entropy and KL divergence]\label{def:cross-entropy}
The \term{cross-entropy} of $\vq$ relative to $\vp$ and the
\term{Kullback--Leibler divergence} (KL divergence) of $\vq$ from $\vp$ are
\[
  H(\vp, \vq) = -\sum_{k=1}^K p_k\log q_k,
  \qquad
  \keyeq{D_{\KL}(\vp \,\Vert\, \vq) = \sum_{k=1}^K p_k\log\frac{p_k}{q_k}
  = H(\vp, \vq) - H(\vp)}
\]
\end{definition}

The cross-entropy is the average code length when we are wrong about the
distribution, and the KL divergence is the price of being wrong: the number of
extra bits (or nats) per outcome.

\begin{theorem}[Gibbs' inequality]\label{thm:gibbs}
$D_{\KL}(\vp \,\Vert\, \vq) \ge 0$, with equality only if $\vp = \vq$. Consequently
$H(\vp, \vq) \ge H(\vp)$.
\end{theorem}

\begin{proof}
We use natural logarithms and the inequality $\ln t \le t - 1$, which says that the
concave function $\ln$ lies below its tangent line at $t = 1$, touching it only
at $t = 1$ (\cref{sec:convexity}). Summing over the $k$ with $p_k > 0$,
\[
  -D_{\KL}(\vp \,\Vert\, \vq) = \sum_kp_k\ln\frac{q_k}{p_k}
  \le \sum_kp_k\left(\frac{q_k}{p_k} - 1\right) = \sum_kq_k - \sum_kp_k
  \le 1 - 1 = 0 .
\]
Equality requires $\frac{q_k}{p_k} = 1$ for all these $k$, i.e.\ $\vq = \vp$.
\end{proof}

\begin{example}[KL divergence is not symmetric]\label{ex:kl}
Let $\vp = (0.5, 0.5)$ and $\vq = (0.9, 0.1)$. In nats,
\begin{align*}
  D_{\KL}(\vp \,\Vert\, \vq) &= 0.5\ln\tfrac{0.5}{0.9} + 0.5\ln\tfrac{0.5}{0.1} \approx 0.511,\\
  D_{\KL}(\vq \,\Vert\, \vp) &= 0.9\ln\tfrac{0.9}{0.5} + 0.1\ln\tfrac{0.1}{0.5} \approx 0.368 .
\end{align*}
The KL divergence measures how different two distributions are, but it is not
a distance in the sense of \cref{sec:distance}: it is not symmetric.
\end{example}

\subsection{Cross-entropy as a loss function}

For one training example of class $c$, the ``true'' distribution is the one-hot
vector $\vy$ (with $y_c = 1$), and the model predicts $\vq = \softmax(\vz)$. Then
\[
  H(\vy, \vq) = -\sum_ky_k\ln q_k = -\ln q_c ,
\]
the negative log-probability of the correct class: precisely the negative
log-likelihood of \cref{sec:losses-from-likelihood}. Since a one-hot vector has
entropy $H(\vy) = 0$, this is also $D_{\KL}(\vy \,\Vert\, \vq)$.

\begin{keytakeaways}[Three names for one loss]
Averaged over a training set, the following are the same thing:
\begin{itemize}
  \item the negative log-likelihood of the data (\cref{sec:mle});
  \item the cross-entropy between the labels and the predictions;
  \item the KL divergence from the labels to the predictions, up to a constant
    that does not depend on the model.
\end{itemize}
Training a classifier by minimising cross-entropy means making its predicted
distribution as close as possible to the observed one.
\end{keytakeaways}

\begin{mlnote}[Information theory in machine learning]
Entropy, cross-entropy and KL divergence turn up all over machine learning.
\begin{itemize}
  \item \textbf{Perplexity.} Language models report $e^{\text{cross-entropy}}$,
    the \term{perplexity}. A model with an average cross-entropy of $3$ nats per
    token has perplexity $e^3 \approx 20$: it is as uncertain as if it were choosing
    uniformly among $20$ tokens. Guessing uniformly from a vocabulary of $50\,000$
    tokens gives a cross-entropy of $\ln50\,000 \approx 10.8$ nats.
  \item \textbf{KL divergence} appears in variational autoencoders, in
    \emph{knowledge distillation} (training a small model to match the softened
    output distribution of a large one), and in fine-tuning language models with
    human feedback, where a KL penalty keeps the new model close to the original.
  \item \textbf{Maximum entropy.} For a continuous distribution with density
    $p$, the entropy is defined as $-\int p(x)\ln p(x)\,\dd x$. Among all
    distributions with a given variance, the normal distribution has the largest
    entropy, which is one more reason why it is the default model of noise
    (\cref{sec:distributions}).
\end{itemize}
\end{mlnote}

\sectionexercises

\begin{exercise}
Compute the entropy in bits of $\left(\tfrac14, \tfrac14, \tfrac14, \tfrac14\right)$
and of $\left(\tfrac12, \tfrac12, 0, 0\right)$.
\end{exercise}
\begin{solution}
$4\cdot\tfrac14\cdot2 = 2$ bits, and $2\cdot\tfrac12\cdot1 = 1$ bit (outcomes of
probability $0$ contribute nothing).
\end{solution}

\begin{exercise}
A model predicts $\vq = (0.7, 0.2, 0.1)$. What is the cross-entropy loss if the
true class is $1$? If it is $2$?
\end{exercise}
\begin{solution}
$-\ln0.7 \approx 0.357$ and $-\ln0.2 \approx 1.609$.
\end{solution}

\begin{exercise}
Compute $D_{\KL}(\vp \,\Vert\, \vq)$ in bits for $\vp = \left(\tfrac12, \tfrac12\right)$
and $\vq = \left(\tfrac14, \tfrac34\right)$.
\end{exercise}
\begin{solution}
$\tfrac12\log_2\tfrac{1/2}{1/4} + \tfrac12\log_2\tfrac{1/2}{3/4} = \tfrac12\cdot1 +
\tfrac12\log_2\tfrac23 \approx 0.5 - 0.292 = 0.208$ bits.
\end{solution}

\begin{exercise}
Use Gibbs' inequality with $\vq$ the uniform distribution on $K$ outcomes to show
that $H(\vp) \le \log K$ for every distribution $\vp$ on $K$ outcomes.
\end{exercise}
\begin{solution}
With $q_k = \tfrac1K$, $D_{\KL}(\vp \,\Vert\, \vq) = \sum_kp_k\log(Kp_k) = \log K
- H(\vp)$, using $\sum_kp_k = 1$. Since this is at least $0$, $H(\vp) \le \log K$,
with equality only for the uniform distribution.
\end{solution}

\begin{exercise}
A language model has an average cross-entropy of $2.3$ nats per token on a test
text. What is its perplexity?
\end{exercise}
\begin{solution}
$e^{2.3} \approx 10$: on average it is as uncertain as a uniform choice among
about $10$ tokens.
\end{solution}

%----------------------------------------------------------------------
\section{Matrix calculus}\label{sec:matrix-calculus}
%----------------------------------------------------------------------

Backpropagation needs derivatives of a scalar loss with respect to vectors and
matrices. We follow one simple convention.

\begin{notation}
For a scalar $L$ and a vector $\vx \in \R^n$, $\pderiv{L}{\vx} = \nabla_{\vx}L$ is the
vector of partial derivatives $\pderiv{L}{x_i}$ (\cref{sec:gradient}). For a matrix
$\mW$, $\pderiv{L}{\mW}$ is the matrix with entries $\pderiv{L}{W_{ij}}$. In both
cases, \emph{the derivative has the same shape as the variable}.
\end{notation}

The reliable method for finding such derivatives is to write everything in
components, differentiate with respect to one entry, and reassemble the result
into vectors and matrices.

\begin{theorem}[Basic gradients]\label{thm:basic-gradients}
For a constant vector $\va$, a constant matrix $\mA$ and a constant vector
$\vb$:
\begin{enumerate}
  \item $\nabla_{\vx}(\va^\top\vx) = \va$;
  \item $\nabla_{\vx}(\vx^\top\mA\vx) = (\mA + \mA^\top)\vx$, which equals
    $2\mA\vx$ if $\mA$ is symmetric;
  \item $\nabla_{\vx}\norm{\vx}^2 = 2\vx$;
  \item $\nabla_{\vx}\norm{\mA\vx - \vb}^2 = 2\mA^\top(\mA\vx - \vb)$.
\end{enumerate}
\end{theorem}

\begin{proof}
(1) $\va^\top\vx = \sum_ia_ix_i$, whose derivative with respect to $x_k$ is $a_k$.
(2) $\vx^\top\mA\vx = \sum_i\sum_jx_iA_{ij}x_j$. The variable $x_k$ appears when
$i = k$ and when $j = k$, so
$\pderiv{}{x_k} = \sum_jA_{kj}x_j + \sum_ix_iA_{ik} = (\mA\vx)_k + (\mA^\top\vx)_k$.
(3) is (2) with $\mA = \mI$. (4) follows from (3) and the chain rule: with
$\vr = \mA\vx - \vb$, \cref{thm:linear-layer-gradients} below gives
$\nabla_{\vx} = \mA^\top\nabla_{\vr}\norm{\vr}^2 = 2\mA^\top\vr$.
\end{proof}

Setting the gradient in (4) to zero gives $\mA^\top\mA\vx = \mA^\top\vb$, the normal
equations of least squares (\cref{thm:normal-equations}), found this time with
calculus instead of geometry.

\begin{example}[Checking an identity]\label{ex:quadratic-gradient}
Let $\mA = \left(\begin{smallmatrix} 1 & 2\\ 3 & 4\end{smallmatrix}\right)$ and
$f(\vx) = \vx^\top\mA\vx = x_1^2 + 5x_1x_2 + 4x_2^2$. At $\vx = (1, 1)$ the formula
gives $(\mA + \mA^\top)\vx = \left(\begin{smallmatrix} 2 & 5\\ 5 &
8\end{smallmatrix}\right)\left(\begin{smallmatrix} 1\\ 1\end{smallmatrix}\right)
= (7, 13)$. Directly: $\pderiv{f}{x_1} = 2x_1 + 5x_2 = 7$ and
$\pderiv{f}{x_2} = 5x_1 + 8x_2 = 13$. \checkmark
\end{example}

\subsection{Gradients through a layer}

The following two rules are all that backpropagation needs.

\begin{theorem}[Gradients through a linear layer]\label{thm:linear-layer-gradients}
Let $\vz = \mW\vx + \vb$, with $\mW$ an $m\times n$ matrix, and let $L$ be a scalar
that depends on $\vz$, with gradient $\vdelta = \pderiv{L}{\vz} \in \R^m$. Then
\[
  \keyeq{\pderiv{L}{\mW} = \vdelta\,\vx^\top, \qquad
  \pderiv{L}{\vx} = \mW^\top\vdelta, \qquad
  \pderiv{L}{\vb} = \vdelta}
\]
\end{theorem}

\begin{proof}
In components, $z_i = \sum_jW_{ij}x_j + b_i$. The weight $W_{ij}$ appears only in
$z_i$, with $\pderiv{z_i}{W_{ij}} = x_j$, so $\pderiv{L}{W_{ij}} = \delta_ix_j$:
the matrix of these numbers is the outer product $\vdelta\vx^\top$
(\cref{sec:transpose}). The input $x_j$ appears in every $z_i$, with
$\pderiv{z_i}{x_j} = W_{ij}$, so by the multivariable chain rule
$\pderiv{L}{x_j} = \sum_i\delta_iW_{ij} = (\mW^\top\vdelta)_j$. Finally,
$\pderiv{z_i}{b_i} = 1$.
\end{proof}

\begin{intuition}[Reading the weight gradient]
The weight $W_{ij}$ connects input $j$ to output $i$, and its gradient is
$\delta_ix_j$: (how much the loss cares about output $i$) times (how active input
$j$ was). A weight whose input was $0$ gets no gradient, since it had no
influence on this example.
\end{intuition}

For the activation step we need a new product. The \term{Hadamard product} of two
vectors of the same length multiplies them component by component:
$\vu \odot \vv = (u_1v_1, \dots, u_nv_n)$.

\begin{theorem}[Gradients through an activation]\label{thm:activation-gradients}
If $\va = \varphi(\vz)$ componentwise and $L$ depends on $\va$, then
\[
  \keyeq{\pderiv{L}{\vz} = \pderiv{L}{\va} \odot \varphi'(\vz)}
\]
\end{theorem}

\begin{proof}
Each $a_i = \varphi(z_i)$ depends only on $z_i$, so
$\pderiv{L}{z_i} = \pderiv{L}{a_i}\varphi'(z_i)$. (In matrix language, the Jacobian
of $\va$ is diagonal, as in \cref{ex:layer-jacobian}.)
\end{proof}

\begin{remark}[Let the shapes guide you]
Shapes catch most mistakes in matrix calculus. If $\mW$ is $m\times n$, then
$\pderiv{L}{\mW}$ must be $m\times n$ too, and the only way to build such a matrix
from $\vdelta \in \R^m$ and $\vx \in \R^n$ is $\vdelta\vx^\top$. Likewise
$\pderiv{L}{\vx} \in \R^n$ can only be $\mW^\top\vdelta$. For a mini-batch stored
as a matrix $\mX$ with one example per row, the forward pass is
$\mZ = \mX\mW^\top + \vone\vb^\top$, and with $\mat{\Delta} = \pderiv{L}{\mZ}$ (one
row per example) the gradients become $\pderiv{L}{\mW} = \mat{\Delta}^\top\mX$,
the sum of the outer products of all the examples, and
$\pderiv{L}{\mX} = \mat{\Delta}\mW$.
\end{remark}

\sectionexercises

\begin{exercise}
Find $\nabla_{\vx}\left(\vx^\top\vx + 3\va^\top\vx\right)$.
\end{exercise}
\begin{solution}
$2\vx + 3\va$, by parts (3) and (1) of \cref{thm:basic-gradients}.
\end{solution}

\begin{exercise}
Show that the gradient of the mean squared error
$L(\vw) = \frac1n\norm{\mX\vw - \vy}^2$ is $\frac2n\mX^\top(\mX\vw - \vy)$.
\end{exercise}
\begin{solution}
Apply part (4) of \cref{thm:basic-gradients} with $\mA = \mX$ and $\vb = \vy$,
and multiply by the constant $\tfrac1n$.
\end{solution}

\begin{exercise}
For a layer $\vz = \mW\vx + \vb$ with $\mW = \left(\begin{smallmatrix} 1 & 0 & 2\\
0 & 1 & -1\end{smallmatrix}\right)$ and $\vx = (3, 0, 1)$, suppose
$\vdelta = \pderiv{L}{\vz} = (1, -2)$. Compute $\pderiv{L}{\mW}$, $\pderiv{L}{\vx}$
and $\pderiv{L}{\vb}$.
\end{exercise}
\begin{solution}
$\pderiv{L}{\mW} = \vdelta\vx^\top = \left(\begin{smallmatrix} 3 & 0 & 1\\ -6 & 0 &
-2\end{smallmatrix}\right)$; $\pderiv{L}{\vx} = \mW^\top\vdelta = (1, -2, 2 + 2) = (1,
-2, 4)$; $\pderiv{L}{\vb} = (1, -2)$.
\end{solution}

\begin{exercise}
A layer has a $128\times784$ weight matrix. What is the shape of the gradient of
the loss with respect to it, and of the vector $\vdelta$?
\end{exercise}
\begin{solution}
The gradient is also $128\times784$ ($100\,352$ numbers); $\vdelta$ has $128$
components, one per neuron.
\end{solution}

%----------------------------------------------------------------------
\section{The gradient of the cross-entropy loss}\label{sec:softmax-gradient}
%----------------------------------------------------------------------

We now differentiate the most important loss in deep learning: softmax followed by
cross-entropy. The result is remarkably simple.

\begin{theorem}[Gradient of softmax with cross-entropy]\label{thm:softmax-ce-gradient}
Let $\vq = \softmax(\vz)$ and $L = -\ln q_c$, where $c$ is the correct class, and
let $\vy$ be the one-hot vector of $c$. Then
\[
  \keyeq{\pderiv{L}{\vz} = \vq - \vy}
\]
\end{theorem}

\begin{proof}
By the definition of softmax,
\[
  L = -\ln\frac{e^{z_c}}{\sum_je^{z_j}} = -z_c + \ln\sum_je^{z_j} .
\]
The derivative of $-z_c$ with respect to $z_k$ is $-1$ if $k = c$ and $0$
otherwise, that is, $-y_k$. By the chain rule, the derivative of
$\ln\sum_je^{z_j}$ is $\frac{e^{z_k}}{\sum_je^{z_j}} = q_k$. Together,
$\pderiv{L}{z_k} = q_k - y_k$.
\end{proof}

\begin{intuition}[Prediction minus truth]
The gradient is simply ``what the model predicted minus what it should have
predicted''. For the correct class, $q_c - 1 < 0$, so gradient descent raises its
logit; every other logit is lowered in proportion to the probability it wrongly
received. A perfect prediction gives a zero gradient, and the gradient can never
be larger than $1$ in any component.
\end{intuition}

\begin{example}[One step on the logits]\label{ex:softmax-step}
In \cref{ex:softmax}, $\vz = (2, 1, 0.1)$ and $\vq \approx (0.659, 0.242, 0.099)$.
If the correct class is the first, the loss is $-\ln0.659 \approx 0.417$ and the
gradient is
\[
  \vq - \vy \approx (-0.341,\ 0.242,\ 0.099).
\]
A gradient step of size $1$ on the logits gives $\vz \approx (2.341, 0.758,
0.001)$, for which $q_1 \approx 0.768$ and the loss has dropped to $0.264$.
\end{example}

The same pattern holds for the other standard pairs of output and loss. For a
sigmoid output $q = \sigma(z)$ with binary cross-entropy, the chain rule and
$\sigma' = \sigma(1 - \sigma)$ (\cref{ex:sigmoid-derivative}) give
\[
  \pderiv{L}{z} = \left(-\frac yq + \frac{1 - y}{1 - q}\right)q(1 - q) = q - y ,
\]
and for a linear output $\hat y = z$ with the squared error $\frac12(\hat y - y)^2$,
$\pderiv{L}{z} = \hat y - y$. In each case: prediction minus target. Combined with
\cref{thm:linear-layer-gradients}, logistic regression ($z = \vw^\top\vx + b$) has
the gradient
\[
  \nabla_{\vw}L = (q - y)\,\vx, \qquad \pderiv{L}{b} = q - y .
\]

\begin{remark}[Why not squared error for classification?]
Suppose we used the squared error $\frac12(q - y)^2$ on a sigmoid output. Its
gradient with respect to the logit is $(q - y)\,q(1 - q)$. When the model is
confidently wrong, say $z = -5$ for an example with $y = 1$, we get
$q = \sigma(-5) \approx 0.0067$ and a gradient of only about $-0.0066$, while the
cross-entropy gives $q - y \approx -0.993$, about $150$ times larger. With squared
error the factor $q(1 - q)$ silences exactly the examples that most need
correcting; cross-entropy cancels it.
\end{remark}

\begin{pitfall}[Compute the loss from the logits]
Never compute $\softmax(\vz)$ and then take its logarithm: a tiny probability can
underflow to $0$, and $-\ln0 = \infty$. Use $-\ln q_c = -z_c + \ln\sum_je^{z_j}$ with
the stable log-sum-exp of \cref{sec:softmax}. Deep learning libraries offer
``cross-entropy with logits'' functions for exactly this reason.
\end{pitfall}

\sectionexercises

\begin{exercise}
The logits are $\vz = (0, 0, 0)$ and the correct class is $3$. Compute the
probabilities, the loss and the gradient.
\end{exercise}
\begin{solution}
$\vq = \left(\tfrac13, \tfrac13, \tfrac13\right)$, $L = \ln3 \approx 1.099$ and
$\vq - \vy = \left(\tfrac13, \tfrac13, -\tfrac23\right)$.
\end{solution}

\begin{exercise}
Show that the components of $\vq - \vy$ always add up to $0$. Explain this using
the shift invariance of softmax.
\end{exercise}
\begin{solution}
$\sum_k(q_k - y_k) = 1 - 1 = 0$. Adding the same constant to all logits does not
change the loss, so the derivative of the loss in the direction $\vone =
(1, \dots, 1)$, which is $\vone^\top(\vq - \vy)$, must be $0$.
\end{solution}

\begin{exercise}
A binary classifier outputs the logit $z = 2$ for an example with label $y = 0$.
Find $q$, the loss and $\pderiv{L}{z}$.
\end{exercise}
\begin{solution}
$q = \sigma(2) \approx 0.881$, $L = -\ln(1 - 0.881) \approx 2.127$ and
$\pderiv{L}{z} = q - y \approx 0.881$.
\end{solution}

\begin{exercise}
Logistic regression has $\vw = (1, -1)$ and $b = 0$. For $\vx = (2, 1)$ and $y = 1$,
compute the gradients of the cross-entropy loss with respect to $\vw$ and $b$.
\end{exercise}
\begin{solution}
$z = 2 - 1 = 1$, $q = \sigma(1) \approx 0.731$, so $\nabla_{\vw}L = (q - 1)\vx \approx
-0.269\,(2, 1) = (-0.538, -0.269)$ and $\pderiv{L}{b} \approx -0.269$.
\end{solution}

%----------------------------------------------------------------------
\section{Backpropagation}\label{sec:backpropagation}
%----------------------------------------------------------------------

We now have everything needed to compute the gradient of the loss with respect to
every weight of a network. Write $\vdelta^{(l)} = \pderiv{L}{\vz^{(l)}}$ for the
gradient with respect to the pre-activations of layer $l$, sometimes called the
\emph{error signal} of that layer.

\begin{theorem}[The backpropagation equations]\label{thm:backprop}
For the network of \cref{def:mlp} and a loss $L$ computed from its output:
\begin{align*}
  \vdelta^{(L)} &= \pderiv{L}{\vz^{(L)}}
    &&\text{(for example $\vq - \vy$, by \cref{thm:softmax-ce-gradient})},\\
  \vdelta^{(l)} &= \bigl(\mW^{(l+1)\top}\vdelta^{(l+1)}\bigr) \odot \varphi'\bigl(\vz^{(l)}\bigr)
    &&\text{for } l = L - 1, \dots, 1,\\
  \pderiv{L}{\mW^{(l)}} &= \vdelta^{(l)}\,\va^{(l-1)\top},
  \qquad \pderiv{L}{\vb^{(l)}} = \vdelta^{(l)}
    &&\text{for } l = 1, \dots, L .
\end{align*}
\end{theorem}

\begin{proof}
Layer $l + 1$ computes $\vz^{(l+1)} = \mW^{(l+1)}\va^{(l)} + \vb^{(l+1)}$, so by
\cref{thm:linear-layer-gradients}, $\pderiv{L}{\va^{(l)}} = \mW^{(l+1)\top}\vdelta^{(l+1)}$;
then $\va^{(l)} = \varphi(\vz^{(l)})$ and \cref{thm:activation-gradients} give the
second equation. The weight and bias gradients follow from
\cref{thm:linear-layer-gradients} applied to layer $l$, whose input is
$\va^{(l-1)}$.
\end{proof}

\begin{algorithm}[Backpropagation]\label{alg:backprop}
To compute the gradient of the loss for one example $(\vx, \vy)$:
\begin{steps}
  \item \textbf{Forward pass.} Set $\va^{(0)} = \vx$ and compute, for
    $l = 1, \dots, L$, the vectors $\vz^{(l)}$ and $\va^{(l)}$, storing them all.
    Compute the loss.
  \item \textbf{Output layer.} Compute $\vdelta^{(L)} = \pderiv{L}{\vz^{(L)}}$.
  \item \textbf{Backward pass.} For $l = L, L - 1, \dots, 1$: record the gradients
    $\vdelta^{(l)}\va^{(l-1)\top}$ and $\vdelta^{(l)}$, and if $l > 1$ compute
    $\vdelta^{(l-1)} = \bigl(\mW^{(l)\top}\vdelta^{(l)}\bigr) \odot
    \varphi'\bigl(\vz^{(l-1)}\bigr)$.
\end{steps}
For a mini-batch, average the gradients of its examples (or use the matrix form
of the remark in \cref{sec:matrix-calculus}).
\end{algorithm}
\index{backpropagation}

\begin{example}[Backpropagation by hand]\label{ex:backprop}
A network has $2$ inputs, a hidden layer of $2$ ReLU neurons and $1$ linear
output, with
\[
  \mW^{(1)} = \begin{pmatrix} 1 & -1\\ 0.5 & 1\end{pmatrix},\quad
  \vb^{(1)} = \begin{pmatrix} 0\\ -1\end{pmatrix},\quad
  \mW^{(2)} = \begin{pmatrix} 2 & -1\end{pmatrix},\quad
  b^{(2)} = 0.5 ,
\]
and the loss $L = \frac12(\hat y - y)^2$. Compute all gradients for the example
$\vx = (1, 2)$, $y = 1$, and take one gradient descent step with $\eta = 0.1$.

\sol \emph{Forward pass.}
$\vz^{(1)} = (1 - 2 + 0,\ 0.5 + 2 - 1) = (-1, 1.5)$, so
$\va^{(1)} = \ReLU(\vz^{(1)}) = (0, 1.5)$. Then
$\hat y = z^{(2)} = 2\cdot0 - 1\cdot1.5 + 0.5 = -1$, and $L = \frac12(-1 - 1)^2 = 2$.

\emph{Output layer.} $\delta^{(2)} = \hat y - y = -2$, so
\[
  \pderiv{L}{\mW^{(2)}} = \delta^{(2)}\va^{(1)\top} = (0,\ -3), \qquad
  \pderiv{L}{b^{(2)}} = -2 .
\]

\emph{Hidden layer.} $\mW^{(2)\top}\delta^{(2)} = (-4, 2)$, and
$\ReLU'(\vz^{(1)}) = (0, 1)$ because the first neuron is inactive. So
$\vdelta^{(1)} = (-4, 2) \odot (0, 1) = (0, 2)$ and
\[
  \pderiv{L}{\mW^{(1)}} = \vdelta^{(1)}\vx^\top
  = \begin{pmatrix} 0\\ 2\end{pmatrix}\begin{pmatrix} 1 & 2\end{pmatrix}
  = \begin{pmatrix} 0 & 0\\ 2 & 4\end{pmatrix}, \qquad
  \pderiv{L}{\vb^{(1)}} = (0, 2).
\]

\emph{Update.} Subtracting $0.1$ times each gradient gives
$\mW^{(1)} = \left(\begin{smallmatrix} 1 & -1\\ 0.3 & 0.6\end{smallmatrix}\right)$,
$\vb^{(1)} = (0, -1.2)$, $\mW^{(2)} = (2\ \ {-0.7})$ and $b^{(2)} = 0.7$. Running the
forward pass again: $\vz^{(1)} = (-1, 0.3)$, $\va^{(1)} = (0, 0.3)$,
$\hat y = -0.21 + 0.7 = 0.49$, and the loss has fallen from $2$ to
$\frac12(0.49 - 1)^2 \approx 0.13$.
\end{example}

\begin{figure}[htbp]
\centering
\begin{tikzpicture}[>={Stealth[length=2.2mm]},
  neuron/.style={circle,draw,thick,minimum size=10mm,inner sep=0pt,font=\small},
  fwd/.style={font=\small,text=PlotBlue},
  bwd/.style={font=\small,text=PlotRed}]
  \node[neuron,fill=MLdef!12] (x1) at (0,1.4) {$x_1$};
  \node[neuron,fill=MLdef!12] (x2) at (0,-1.4) {$x_2$};
  \node[neuron,fill=MLexample!14] (h1) at (4,1.4) {$h_1$};
  \node[neuron,fill=MLexample!14] (h2) at (4,-1.4) {$h_2$};
  \node[neuron,fill=MLintuition!22] (o) at (8,0) {$\hat y$};
  \draw[->,thick,gray] (x1) -- node[pos=0.22,above,font=\scriptsize,text=black] {$1$} (h1);
  \draw[->,thick,gray] (x2) -- node[pos=0.22,below,font=\scriptsize,text=black] {$1$} (h2);
  \draw[->,thick,gray] (x1) -- node[pos=0.2,above right=-1pt,font=\scriptsize,text=black] {$0.5$} (h2);
  \draw[->,thick,gray] (x2) -- node[pos=0.2,below right=-1pt,font=\scriptsize,text=black] {$-1$} (h1);
  \draw[->,thick,gray] (h1) -- node[above,font=\scriptsize,text=black] {$2$} (o);
  \draw[->,thick,gray] (h2) -- node[below,font=\scriptsize,text=black] {$-1$} (o);
  \node[fwd,left=3pt] at (x1.west) {$1$};
  \node[fwd,left=3pt] at (x2.west) {$2$};
  \node[fwd,above=2pt,align=center] at (h1.north) {$z = -1$, $a = 0$};
  \node[fwd,below=2pt,align=center] at (h2.south) {$z = 1.5$, $a = 1.5$};
  \node[fwd,above=2pt] at (o.north) {$\hat y = -1$, $L = 2$};
  \node[bwd,right=3pt] at (h1.east) {$\delta = 0$};
  \node[bwd,right=3pt] at (h2.east) {$\delta = 2$};
  \node[bwd,below=2pt] at (o.south) {$\delta = -2$};
\end{tikzpicture}
\caption{The network of \cref{ex:backprop}, with the weights in black, the
  forward values in blue and the error signals $\delta = \pderiv{L}{z}$ in red.
  The inactive neuron $h_1$ passes no gradient back.}
\label{fig:backprop}
\end{figure}

\begin{intuition}[Why backpropagation is fast]
The backward pass costs about as much as the forward pass: one multiplication by
$\mW^{(l)\top}$ per layer, plus an outer product. So all the partial derivatives,
for millions or billions of weights, cost only a few forward passes. Numerical
differentiation (\cref{sec:numerical-derivatives}) would need a separate forward
pass for every single weight. The price is memory: every $\vz^{(l)}$ and
$\va^{(l)}$ of the forward pass must be stored until the backward pass uses it,
which is why training needs much more memory than using a trained network.
\end{intuition}

\sectionexercises

\begin{exercise}
In \cref{ex:backprop}, what are the gradients if the target is $y = -1$ instead
of $1$?
\end{exercise}
\begin{solution}
Then $\hat y = y$, the loss is $0$, $\delta^{(2)} = 0$, and all gradients are $0$:
a perfectly predicted example teaches the network nothing.
\end{solution}

\begin{exercise}
Why is the first row of $\pderiv{L}{\mW^{(1)}}$ zero in \cref{ex:backprop}? What
would happen if that neuron were inactive for every training example?
\end{exercise}
\begin{solution}
The first hidden neuron has $z = -1 < 0$, so $\ReLU'(z) = 0$ and $\delta_1^{(1)} = 0$:
its weights had no effect on the output for this input. If it were inactive for
all examples, its weights would never change: a ``dead'' ReLU neuron.
\end{solution}

\begin{exercise}
A network has one input, one hidden ReLU neuron and one linear output:
$h = \ReLU(w_1x + b_1)$, $\hat y = w_2h + b_2$, with loss $\frac12(\hat y - y)^2$. For
$x = 1$, $w_1 = 2$, $b_1 = -1$, $w_2 = 3$, $b_2 = 0$ and $y = 2$, compute the gradient
with respect to all four parameters.
\end{exercise}
\begin{solution}
Forward: $z_1 = 1$, $h = 1$, $\hat y = 3$, $L = \tfrac12$. Backward: $\delta_2 = \hat y -
y = 1$, so $\pderiv{L}{w_2} = \delta_2h = 1$ and $\pderiv{L}{b_2} = 1$. Then
$\delta_1 = w_2\delta_2\ReLU'(z_1) = 3$, so $\pderiv{L}{w_1} = \delta_1x = 3$ and
$\pderiv{L}{b_1} = 3$.
\end{solution}

\begin{exercise}
A network has layer sizes $3 \to 4 \to 2$. Give the shapes of $\vdelta^{(1)}$,
$\vdelta^{(2)}$, $\pderiv{L}{\mW^{(1)}}$ and $\pderiv{L}{\mW^{(2)}}$.
\end{exercise}
\begin{solution}
$\vdelta^{(1)} \in \R^4$ and $\vdelta^{(2)} \in \R^2$; $\pderiv{L}{\mW^{(1)}}$ is
$4\times3$ and $\pderiv{L}{\mW^{(2)}}$ is $2\times4$, the shapes of the weight
matrices.
\end{solution}

%----------------------------------------------------------------------
\section{Vanishing and exploding gradients}\label{sec:vanishing-gradients}
%----------------------------------------------------------------------

The backpropagation equations multiply the error signal by a matrix and by the
activation derivatives once per layer. In a deep network this means long
products, and long products of numbers tend to vanish or explode: in
\cref{sec:exponential} we saw that $0.9^{100} \approx 0.000027$ and
$1.1^{100} \approx 13\,781$.

\begin{picturethis}[A chain of one-neuron layers]
Imagine a network with a single neuron per layer, $z^{(l)} = w_la^{(l-1)}$ and
$a^{(l)} = \varphi(z^{(l)})$. By the chain rule,
\[
  \pderiv{L}{a^{(0)}} = \pderiv{L}{a^{(L)}}\prod_{l=1}^L w_l\,\varphi'\bigl(z^{(l)}\bigr).
\]
If every factor $w_l\varphi'(z^{(l)})$ is about $c$, the gradient reaching the first
layer is about $c^L$ times the gradient at the last (\cref{fig:vanishing}). With
$c = 0.75$ and $50$ layers, it shrinks by a factor of about $10^{-6}$; with
$c = 1.25$, it grows by a factor of $70\,000$.
\end{picturethis}

\begin{figure}[htbp]
\centering
\begin{tikzpicture}
\begin{axis}[mlplot,axis lines=left,ymode=log,width=10cm,height=6.2cm,xmin=0,xmax=50,
  ymin=1e-7,ymax=1e5,xlabel={number of layers $L$},ylabel={gradient factor $c^L$},
  xlabel style={at={(ticklabel cs:0.5)},anchor=north},
  ylabel style={at={(ticklabel cs:0.5)},anchor=south,rotate=90},clip=false]
  \addplot[PlotRed,domain=0:50,samples=60]{1.25^x};
  \addplot[PlotOrange,domain=0:50,samples=60]{1.05^x};
  \addplot[gray,dashed,domain=0:50,samples=2]{1};
  \addplot[PlotGreen,domain=0:50,samples=60]{0.95^x};
  \addplot[PlotBlue,domain=0:50,samples=60]{0.75^x};
  \node[font=\small,anchor=west,PlotRed] at (axis cs:50,70065) {$c = 1.25$};
  \node[font=\small,anchor=west,PlotOrange] at (axis cs:50,11.47) {$c = 1.05$};
  \node[font=\small,anchor=west,gray] at (axis cs:50,1) {$c = 1$};
  \node[font=\small,anchor=west,PlotGreen] at (axis cs:50,0.077) {$c = 0.95$};
  \node[font=\small,anchor=west,PlotBlue] at (axis cs:50,5.7e-7) {$c = 0.75$};
\end{axis}
\end{tikzpicture}
\caption{If each layer multiplies the gradient by a factor $c$, the gradient
  reaching the first of $L$ layers is multiplied by $c^L$ (logarithmic scale).
  Only factors very close to $1$ keep it at a usable size.}
\label{fig:vanishing}
\end{figure}

For full layers the factors are matrices, and their effect on the length of
$\vdelta$ is governed by their singular values (\cref{sec:svd}), but the
conclusion is the same: gradients either \term[vanishing gradients]{vanish},
so that the early layers learn extremely slowly, or \term[exploding
gradients]{explode}, so that training becomes unstable.

\begin{figure}[htbp]
\centering
\begin{tikzpicture}
\begin{groupplot}[group style={group size=3 by 1,horizontal sep=0.9cm},
  mlplot,axis lines=left,width=5.2cm,height=4.2cm,xmin=-4,xmax=4,ymin=-0.1,ymax=1.2,
  xtick={-4,-2,0,2,4},ytick={0,0.25,0.5,1},yticklabel style={/pgf/number format/fixed},
  title style={font=\small}]
  \nextgroupplot[title={sigmoid: $\sigma'(z)$}]
    \addplot[PlotBlue,domain=-4:4,samples=100]{exp(-x)/(1+exp(-x))^2};
  \nextgroupplot[title={$\tanh'(z)$}]
    \addplot[PlotBlue,domain=-4:4,samples=100]{1-tanh(x)^2};
  \nextgroupplot[title={$\ReLU'(z)$}]
    \addplot[PlotBlue,domain=-4:0,samples=2]{0};
    \addplot[PlotBlue,domain=0:4,samples=2]{1};
    \addplot[only marks,mark=*,mark size=1.6pt,PlotBlue] coordinates {(0,0)};
    \addplot[only marks,mark=o,mark size=1.6pt,PlotBlue,fill=white] coordinates {(0,1)};
\end{groupplot}
\end{tikzpicture}
\caption{Derivatives of three activation functions. The sigmoid's derivative
  never exceeds $\frac14$; $\tanh'$ reaches $1$ at $0$; the ReLU passes gradients
  through unchanged wherever the neuron is active.}
\label{fig:activation-derivatives}
\end{figure}

\begin{example}[Why deep sigmoid networks were hard to train]\label{ex:sigmoid-vanishing}
The derivative of the sigmoid is $\sigma(1 - \sigma) \le \frac14$, with the maximum at
$z = 0$ (\cref{fig:activation-derivatives}). With weights of size about $1$, every
sigmoid layer shrinks the gradient by a factor of at least $4$; after $10$ layers
the factor is at most $\left(\frac14\right)^{10} \approx 10^{-6}$. The early layers of
such a network barely learn, which is a main reason why deep networks were long
considered untrainable.
\end{example}

\subsection{Remedies}

Modern deep learning combines several defences, each built on mathematics from
this book.
\begin{itemize}
  \item \textbf{ReLU-type activations.} $\ReLU'(z) = 1$ for active neurons, so
    the activations do not shrink the gradient.
  \item \textbf{Careful initialisation.} Choosing the variance of the initial
    weights as $\frac1n$ or $\frac2n$ (Xavier or He initialisation,
    \cref{sec:variance}) keeps the variances of activations and gradients roughly
    constant from layer to layer, so that the factors start out close to $1$.
  \item \textbf{Normalisation layers.} Batch and layer normalisation standardise
    the activations inside the network (\cref{sec:variance}).
  \item \textbf{Residual connections.} A residual block computes
    $\va^{(l+1)} = \va^{(l)} + F(\va^{(l)})$, whose Jacobian is $\mI + \mat{J}_F$.
    The identity term gives the gradient a direct path backwards that does not
    shrink. Residual connections made networks with hundreds of layers trainable
    and are part of every transformer.
  \item \textbf{Gradient clipping.} Against exploding gradients: if
    $\norm{\vg} > \tau$ for a threshold $\tau$, replace $\vg$ by
    $\tau\,\vg/\norm{\vg}$, which keeps the direction but limits the length
    (\cref{sec:vector-length}).
\end{itemize}

\sectionexercises

\begin{exercise}
Bound the factor by which $20$ sigmoid layers with weights of size $1$ can shrink a
gradient.
\end{exercise}
\begin{solution}
Each layer contributes a factor of at most $\tfrac14$, so the total is at most
$\left(\tfrac14\right)^{20} \approx 9.1\times10^{-13}$.
\end{solution}

\begin{exercise}
Each of $20$ blocks applies a function with derivative $f'(x) = -0.05$ (take it to
be constant). Compare the derivative of the composition of the $20$ blocks with
and without residual connections, $x \mapsto x + f(x)$.
\end{exercise}
\begin{solution}
Without residual connections the derivative is $(-0.05)^{20} \approx 9.5\times
10^{-27}$. With them each block has derivative $1 + f'(x) = 0.95$, and the
composition has $0.95^{20} \approx 0.36$.
\end{solution}

\begin{exercise}
Clip the gradient $\vg = (6, -8)$ to length at most $\tau = 5$.
\end{exercise}
\begin{solution}
$\norm{\vg} = 10 > 5$, so $\vg \leftarrow 5\cdot\tfrac{(6, -8)}{10} = (3, -4)$.
\end{solution}

\begin{exercise}
Why do $\tanh$ networks suffer less from vanishing gradients than sigmoid networks?
\end{exercise}
\begin{solution}
$\tanh'(0) = 1$, whereas $\sigma'(0) = \tfrac14$: near zero, $\tanh$ passes gradients
on almost unchanged, while every sigmoid layer divides them by at least $4$.
\end{solution}

%----------------------------------------------------------------------
\section{Automatic differentiation}\label{sec:autodiff}
%----------------------------------------------------------------------

Nobody derives backpropagation formulas by hand for real networks. Libraries
such as PyTorch, JAX and TensorFlow compute gradients by \term{automatic
differentiation} (autodiff), which is neither symbolic differentiation
(manipulating formulas, which grow out of control) nor numerical differentiation
(which is inexact and slow, \cref{sec:numerical-derivatives}).

The idea is simple. Every program, however complicated, is built from elementary
operations ($+$, $\times$, $\exp$, $\ReLU$, matrix products, \dots), and the
derivative of each elementary operation is known. During the forward pass the
library records every operation in a computation graph (\cref{fig:computation-graph}),
together with its local derivatives. The backward pass then visits the graph in
reverse order and applies the chain rule, adding up the contributions of all
paths, exactly as we did by hand in \cref{ex:computation-graph}.

The following complete, working implementation of this idea takes about
forty lines of Python. Each \texttt{Value} remembers the values it was computed
from and the local derivative with respect to each of them.
\begin{lstlisting}
import math

class Value:
    """A number that remembers how it was computed."""

    def __init__(self, data, parents=(), local_grads=()):
        self.data = data                # value from the forward pass
        self.grad = 0.0                 # d(output)/d(self)
        self.parents = parents          # Values it was computed from
        self.local_grads = local_grads  # d(self)/d(parent) for each

    def __add__(self, other):
        other = other if isinstance(other, Value) else Value(other)
        return Value(self.data + other.data, (self, other), (1.0, 1.0))

    def __mul__(self, other):
        other = other if isinstance(other, Value) else Value(other)
        return Value(self.data * other.data, (self, other),
                     (other.data, self.data))

    def relu(self):
        return Value(max(self.data, 0.0), (self,),
                     (1.0 if self.data > 0 else 0.0,))

    def backward(self):
        # List the nodes so that each comes after its parents.
        order, seen = [], set()
        def visit(v):
            if id(v) not in seen:
                seen.add(id(v))
                for p in v.parents:
                    visit(p)
                order.append(v)
        visit(self)
        # Chain rule, from the output back towards the inputs.
        self.grad = 1.0
        for v in reversed(order):
            for p, g in zip(v.parents, v.local_grads):
                p.grad += v.grad * g

x, y = Value(2.0), Value(3.0)
f = (x + y) * (y + 1)
f.backward()
print(f.data, x.grad, y.grad)    # 20.0 4.0 9.0
\end{lstlisting}
The line \texttt{p.grad += v.grad * g} is the chain rule: the derivative of the
output with respect to a parent receives, from each child, the child's gradient
times the local derivative. The \texttt{+=} adds the contributions of different
paths, which is why $y$, used twice, receives $4 + 5 = 9$. The output reproduces
\cref{ex:computation-graph}. Real libraries work the same way, with vectors and
matrices instead of single numbers:
\par\noindent\begin{minipage}{\linewidth}
\begin{lstlisting}
import torch

x = torch.tensor(2.0, requires_grad=True)
y = torch.tensor(3.0, requires_grad=True)
f = (x + y) * (y + 1)
f.backward()
print(x.grad, y.grad)            # tensor(4.) tensor(9.)
\end{lstlisting}
\end{minipage}\par

\begin{remark}[Reverse mode and forward mode]
What we have described is \emph{reverse-mode} automatic differentiation: one
backward pass gives the derivatives of one output with respect to all inputs.
That is exactly the situation in training, with one scalar loss and millions of
parameters. \emph{Forward-mode} autodiff instead carries derivatives along with the
values in the forward pass, giving the derivatives of all outputs with respect to
one input per pass; it is efficient when there are few inputs and many outputs.
\end{remark}

\begin{mlnote}[Gradient checking]
When you implement a new operation or a backward pass by hand, compare its
gradients with the central differences
\[
  \frac{L(\theta + h) - L(\theta - h)}{2h}
\]
of \cref{sec:numerical-derivatives}, on a few random inputs. Relative differences of
about $10^{-7}$ indicate a correct implementation; differences of $10^{-2}$ indicate a
bug.
\end{mlnote}

\sectionexercises

\begin{exercise}
Using the \texttt{Value} class, the program \texttt{x = Value(2.0); f = x * x * x;
f.backward()} computes $f = x^3$. Trace the backward pass and confirm that
\texttt{x.grad} ends up as $12$.
\end{exercise}
\begin{solution}
The forward pass creates $u = x\cdot x = 4$ (local derivatives $2$ and $2$) and
$f = u\cdot x = 8$ (local derivatives $x = 2$ with respect to $u$ and $u = 4$ with
respect to $x$). Backward: $f$ gets gradient $1$; $u$ gets $1\cdot2 = 2$ and $x$ gets
$1\cdot4 = 4$; then $u$ passes $2\cdot2 = 4$ to $x$ twice. In total $4 + 4 + 4 = 12 =
3x^2$.
\end{solution}

\begin{exercise}
Describe how to add a $\tanh$ operation to the \texttt{Value} class.
\end{exercise}
\begin{solution}
Compute $t = \tanh(\text{data})$ and return a new \texttt{Value} with data $t$, the
single parent \texttt{self}, and local derivative $1 - t^2$, since
$\tanh'(x) = 1 - \tanh^2(x)$.
\end{solution}

\begin{exercise}
A model has $10^9$ parameters and one scalar loss. Roughly how many evaluations of
the loss would numerical differentiation need to obtain the full gradient, and how
many passes does reverse-mode autodiff need?
\end{exercise}
\begin{solution}
Central differences need two evaluations per parameter: about $2\times10^9$
evaluations. Reverse mode needs one forward and one backward pass, together
costing a small multiple of a single evaluation.
\end{solution}

%----------------------------------------------------------------------
\section{A glimpse of attention}\label{sec:attention}
%----------------------------------------------------------------------

The transformer architecture behind today's language models is built around one
operation, \emph{attention}. A text is split into tokens, each represented by a
vector; attention lets each token collect information from the others, weighting
them by relevance. Mathematically it uses nothing beyond this book.

\begin{definition}[Scaled dot-product attention]\label{def:attention}
Given a \emph{query} vector $\vq \in \R^d$, \emph{key} vectors
$\vk_1, \dots, \vk_n \in \R^d$ and \emph{value} vectors $\vv_1, \dots, \vv_n$,
attention computes
\[
  s_j = \frac{\vq^\top\vk_j}{\sqrt d}, \qquad
  (\alpha_1, \dots, \alpha_n) = \softmax(s_1, \dots, s_n), \qquad
  \text{output} = \sum_{j=1}^n\alpha_j\vv_j .
\]
For many queries at once, stored as the rows of matrices $\mQ$, $\mK$ and
$\mV$,
\[
  \keyeq{\operatorname{Attention}(\mQ, \mK, \mV)
  = \softmax\!\left(\frac{\mQ\mK^\top}{\sqrt d}\right)\mV}
\]
with the softmax applied to each row.
\end{definition}
\index{attention}

Read it piece by piece: the dot products $\vq^\top\vk_j$ measure how well each key
matches the query (\cref{sec:dot-product}); softmax turns the scores into weights
that are positive and add up to $1$; and the output is the weighted average of the
values, that is, the expected value under the distribution $\alpha$
(\cref{sec:expectation}). In a transformer, the queries, keys and values are
computed from the token vectors by learned matrices, $\mQ = \mX\mW_Q$,
$\mK = \mX\mW_K$ and $\mV = \mX\mW_V$.

\begin{example}[Attention by hand]\label{ex:attention}
Let $d = 2$, keys $\vk_1 = (1, 0)$, $\vk_2 = (0, 1)$, $\vk_3 = (1, 1)$, values equal
to the keys, and the query $\vq = (1, 0)$. The scores are
$\frac{1}{\sqrt2}(1, 0, 1) \approx (0.707, 0, 0.707)$, the weights are
$\softmax(0.707, 0, 0.707) \approx (0.401, 0.198, 0.401)$, and the output is
\[
  0.401\,(1, 0) + 0.198\,(0, 1) + 0.401\,(1, 1) \approx (0.802, 0.599).
\]
The query attends mostly to the two keys that point in its direction.
\end{example}

\begin{intuition}[Why divide by $\sqrt d$?]
If the components of $\vq$ and $\vk$ are independent with mean $0$ and variance $1$,
then $\vq^\top\vk = \sum_{i=1}^dq_ik_i$ is a sum of $d$ independent terms of variance
$1$, so it has variance $d$ (\cref{sec:variance}). For $d = 64$ the scores would
have standard deviation $8$, and softmax of such large numbers is nearly one-hot,
with almost zero gradients (\cref{sec:softmax-gradient}). Dividing by $\sqrt d$
brings the variance back to $1$.
\end{intuition}

\sectionexercises

\begin{exercise}
Repeat \cref{ex:attention} with the query $\vq = (0, 1)$.
\end{exercise}
\begin{solution}
Scores $\tfrac{1}{\sqrt2}(0, 1, 1) \approx (0, 0.707, 0.707)$, weights
$\approx (0.198, 0.401, 0.401)$, output $\approx 0.198(1, 0) + 0.401(0, 1) + 0.401(1,
1) = (0.599, 0.802)$.
\end{solution}

\begin{exercise}
What is the output of attention if all the scores are equal? If one score is much
larger than all the others?
\end{exercise}
\begin{solution}
Equal scores give equal weights $\tfrac1n$, so the output is the plain average of
the values. A dominant score gives a weight close to $1$ for that key, so the output
is close to its value: attention then acts like a look-up.
\end{solution}

\begin{exercise}
Queries and keys have $d = 256$ components, each with mean $0$ and variance $1$,
independently. What is the standard deviation of $\vq^\top\vk$, and of the scaled
score?
\end{exercise}
\begin{solution}
The variance is $256$, so the standard deviation is $16$; after dividing by
$\sqrt{256} = 16$ it is $1$.
\end{solution}

%----------------------------------------------------------------------
\section{Putting it all together: training a network}\label{sec:training-loop}
%----------------------------------------------------------------------

We can now describe, and justify step by step, the complete procedure by which
neural networks are trained.

\begin{algorithm}[Training a neural network]\label{alg:training}
To train a classifier:
\begin{steps}
  \item \textbf{Prepare the data.} Split it into training, validation and test
    sets, and standardise the features using statistics of the training set
    (\cref{sec:generalisation,sec:variance}).
  \item \textbf{Initialise.} Draw the weights at random with variance $\frac2n$
    or $\frac1n$; set the biases to $0$ (\cref{sec:variance}).
  \item \textbf{Loop over epochs and mini-batches.} For each mini-batch: run the
    forward pass and compute the average cross-entropy loss from the logits
    (\cref{sec:softmax-gradient}); run the backward pass to get the gradients
    (\cref{sec:backpropagation}); update the parameters with SGD, momentum or Adam,
    possibly with weight decay and a learning-rate schedule
    (\cref{ch:optimization,sec:map}).
  \item \textbf{Monitor.} Track the training and validation losses; choose
    hyperparameters and the stopping time on the validation set
    (\cref{sec:generalisation}).
  \item \textbf{Evaluate.} Measure the final model once on the test set, with
    error bars (\cref{sec:clt}).
\end{steps}
\end{algorithm}

The following program carries this out for a small problem: points inside the
unit circle belong to class $1$ and points in a surrounding ring to class $0$
(compare \cref{fig:circle-classes}). No straight line separates the classes, but a
network with one hidden layer of $16$ ReLU neurons learns to. For brevity the
program uses all $400$ training points in every step (full-batch gradient
descent) and no validation set; every line corresponds to an equation of this
chapter.
\begin{lstlisting}
import numpy as np

rng = np.random.default_rng(0)

# Data: class 1 inside the unit circle, class 0 in a ring around it.
n = 400
r = np.concatenate([rng.uniform(0, 1, n // 2),
                    rng.uniform(1.5, 2.5, n // 2)])
angle = rng.uniform(0, 2 * np.pi, n)
X = np.column_stack([r * np.cos(angle), r * np.sin(angle)])
y = np.concatenate([np.ones(n // 2), np.zeros(n // 2)])

# One hidden layer of 16 ReLU units and one output logit.
H = 16
W1 = rng.normal(0, np.sqrt(2 / 2), (H, 2))   # He initialisation
b1 = np.zeros(H)
w2 = rng.normal(0, np.sqrt(1 / H), H)
b2 = 0.0

lr = 0.5
for step in range(1, 501):
    # Forward pass for all examples at once (one row each).
    Z1 = X @ W1.T + b1              # pre-activations, n x H
    A1 = np.maximum(Z1, 0)          # ReLU
    z2 = A1 @ w2 + b2               # output logits, length n
    # Binary cross-entropy, computed stably from the logits.
    loss = np.mean(np.logaddexp(0, z2) - y * z2)

    # Backward pass.
    q = 1 / (1 + np.exp(-z2))       # predicted P(y = 1)
    d2 = (q - y) / n                # dL/dz2
    grad_w2 = A1.T @ d2
    grad_b2 = d2.sum()
    D1 = np.outer(d2, w2) * (Z1 > 0)    # dL/dZ1
    grad_W1 = D1.T @ X
    grad_b1 = D1.sum(axis=0)

    # Gradient descent step.
    W1 -= lr * grad_W1
    b1 -= lr * grad_b1
    w2 -= lr * grad_w2
    b2 -= lr * grad_b2

    if step in (1, 10, 50, 100, 500):
        accuracy = np.mean((z2 > 0) == y)
        print(f"step {step:3d}  loss {loss:.3f}  accuracy {accuracy:.3f}")
\end{lstlisting}
It prints:
\par\noindent\begin{minipage}{\linewidth}
\begin{lstlisting}[language={},keywordstyle=\color{black}]
step   1  loss 0.815  accuracy 0.552
step  10  loss 0.445  accuracy 0.807
step  50  loss 0.165  accuracy 1.000
step 100  loss 0.073  accuracy 1.000
step 500  loss 0.010  accuracy 1.000
\end{lstlisting}
\end{minipage}\par
The loss falls steadily, and after $50$ steps every training point is classified
correctly. On $1000$ new points from the same distribution the network is also
right every time: it generalises, because the gap between the two classes leaves
room for a decision boundary that is not tuned to individual training points
(\cref{fig:training-result}).

\begin{figure}[htbp]
\centering
\begin{tikzpicture}
\begin{groupplot}[group style={group size=2 by 1,horizontal sep=1.8cm},height=6cm]
  \nextgroupplot[mlplot,axis lines=box,width=6cm,xmin=-2.7,xmax=2.7,ymin=-2.7,ymax=2.7,
    axis equal image,xtick={-2,0,2},ytick={-2,0,2},grid=none,
    title={data and learned boundary},title style={font=\small}]
    \addplot[only marks,mark=*,mark size=0.9pt,PlotBlue] coordinates {(0.19,0.61) (0.25,-0.10) (0.03,0.02) (0.02,0.00) (-0.36,0.73) (0.91,-0.05) (-0.06,0.60)
      (0.35,-0.64) (0.25,0.48) (-0.80,-0.48) (0.79,-0.21) (-0.00,-0.00) (0.85,-0.10)
      (-0.03,-0.02) (0.73,-0.08) (0.09,-0.15) (0.15,-0.85) (0.41,-0.35) (-0.20,-0.22)
      (-0.26,0.33) (-0.03,-0.01) (0.02,0.12) (0.12,-0.66) (0.31,0.57) (-0.54,-0.29)
      (-0.37,-0.09) (-0.47,-0.88) (0.06,-0.98) (0.53,0.44) (-0.46,-0.46) (-0.59,0.35)
      (-0.29,-0.26) (-0.05,-0.13) (-0.62,-0.37) (-0.06,-0.52) (-0.31,-0.04) (-0.47,0.11)
      (-0.20,0.87) (0.12,0.93) (-0.12,-0.34) (-0.19,-0.54) (0.11,0.30) (0.59,-0.10)
      (-0.16,-0.30) (-0.38,-0.08) (0.48,-0.75) (-0.23,0.02) (-0.62,0.09) (-0.00,0.08)
      (0.46,0.69) (-0.19,-0.76) (0.13,-0.20) (-0.38,-0.79) (-0.04,0.04) (-0.30,-0.15)
      (-0.14,-0.06) (0.42,-0.17) (-0.61,0.52) (0.12,0.20) (0.04,-0.04) (0.32,-0.25)
      (0.19,0.06) (0.03,0.09) (-0.38,-0.44) (0.07,-0.29) (-0.53,-0.42) (0.07,0.19) (0.70,0.63)
      (-0.36,-0.01) (0.04,-0.10) (0.13,0.62) (0.83,0.42) (-0.42,-0.14) (0.34,0.89) (0.46,0.21)
      (0.06,-0.42) (0.27,-0.56) (-0.80,0.59) (-0.26,0.91) (-0.08,0.45) (-0.49,0.58)
      (-0.44,-0.23) (-0.52,-0.09) (-0.48,0.62) (-0.27,-0.32) (-0.33,-0.66) (-0.66,-0.25)
      (-0.71,0.61) (-0.08,-0.08) (-0.61,-0.40) (-0.50,0.78) (-0.32,0.91) (-0.01,-0.00)
      (-0.66,-0.56) (-0.75,-0.63) (-0.71,0.64) (-0.14,-0.06) (0.97,-0.09) (-0.80,0.39)
      (0.45,-0.69) (0.42,0.23) (0.16,-0.16) (0.75,-0.29) (-0.07,0.92) (0.27,0.02) (-0.54,0.06)
      (0.18,0.40) (0.92,-0.17) (0.03,-0.02) (0.71,-0.18) (-0.49,-0.37) (-0.03,-0.00)
      (0.36,-0.62) (-0.01,-0.01) (0.01,0.76) (0.47,-0.21) (-0.86,0.34) (0.01,-0.07)
      (-0.84,-0.00) (0.03,0.06) (-0.10,0.33) (-0.38,-0.19) (0.60,0.76) (0.56,0.05)
      (-0.24,0.10) (0.02,-0.24) (-0.67,-0.58) (-0.10,0.20) (-0.03,-0.12) (-0.29,0.03)
      (0.59,-0.00) (0.09,-0.55) (0.39,-0.71) (-0.03,0.56) (0.17,0.24) (0.13,0.39) (-0.75,0.34)
      (-0.62,-0.05) (0.33,0.90) (0.07,-0.36) (0.37,-0.41) (-0.24,0.54) (-0.85,-0.04)
      (-0.12,-0.08) (-0.07,-0.40) (0.55,0.73) (-0.01,0.04) (-0.10,-0.82) (-0.38,-0.17)
      (0.67,-0.49) (-0.01,0.00) (-0.30,0.20) (-0.03,0.07) (0.08,0.65) (-0.16,-0.22)
      (-0.06,0.70) (0.61,-0.72) (-0.02,0.13) (-0.40,-0.77) (-0.05,-0.02) (-0.26,-0.28)
      (0.34,-0.26) (0.24,0.43) (0.57,0.79) (0.56,0.54) (0.27,0.14) (-0.26,-0.06) (0.44,0.74)
      (0.31,-0.83) (0.51,0.07) (-0.24,0.24) (-0.98,0.17) (0.07,0.31) (-0.11,0.14) (0.15,0.87)
      (-0.16,0.80) (0.60,-0.30) (-0.83,0.48) (-0.70,0.61) (-0.57,-0.48) (-0.44,-0.74)
      (-0.13,-0.21) (0.12,0.07) (-0.58,-0.33) (-0.06,-0.71) (0.05,-0.16) (-0.34,-0.21)
      (0.62,0.67) (0.49,0.28) (-0.26,0.52) (0.17,-0.09) (-0.52,0.09) (0.41,-0.32) (-0.09,0.02)
      (0.03,-0.98) (-0.57,0.05) (-0.00,-0.01) (-0.32,0.70) (0.75,-0.62) (-0.06,0.59)};
    \addplot[only marks,mark=*,mark size=0.9pt,PlotRed] coordinates {(1.82,0.07) (-0.30,-1.66) (-0.97,-1.95) (-0.94,-1.41) (-0.80,-1.92) (-1.80,-1.08)
      (1.84,1.63) (-0.76,-1.37) (2.00,0.08) (0.92,2.05) (-1.47,0.80) (-1.36,1.31) (1.15,1.06)
      (-2.00,0.99) (-1.13,-1.11) (-1.36,1.32) (-0.61,-2.29) (0.24,1.96) (1.49,1.90)
      (-0.02,-2.27) (-1.18,-2.11) (-1.47,0.70) (1.03,1.19) (-0.81,-1.34) (-0.00,-2.37)
      (1.10,1.83) (-0.74,-1.85) (-1.02,1.31) (1.87,-1.11) (0.02,-1.82) (-0.33,2.19)
      (1.81,-0.74) (1.98,0.32) (0.91,2.10) (0.08,1.59) (-0.23,-2.07) (-1.67,-0.28)
      (-2.25,0.51) (0.34,1.96) (0.10,-2.49) (1.25,1.13) (0.04,2.46) (0.80,-2.16) (-1.89,0.60)
      (1.65,-1.62) (-1.69,-1.25) (0.53,-2.09) (0.93,2.23) (-0.63,1.43) (-1.67,1.63)
      (-1.88,0.07) (-1.80,0.31) (1.10,-2.24) (1.06,2.02) (1.18,-1.60) (1.47,-1.23) (2.11,1.09)
      (1.58,0.09) (-0.59,2.13) (-1.86,1.34) (2.26,-0.42) (1.64,0.79) (0.45,-2.25) (-1.70,0.27)
      (1.28,1.36) (-1.28,1.43) (-1.50,1.39) (0.07,1.61) (-0.52,1.83) (-1.31,0.72) (2.18,-0.53)
      (-2.27,0.60) (1.56,-0.51) (2.16,0.42) (2.12,0.94) (2.44,0.43) (-1.18,-2.02) (0.36,1.89)
      (-2.20,-1.15) (0.70,-2.37) (-0.99,1.74) (0.06,2.25) (-0.37,-2.39) (-1.95,0.30)
      (1.40,1.91) (1.88,1.15) (-0.14,-1.79) (1.45,-1.74) (1.60,-1.32) (-1.59,-0.10)
      (1.08,1.55) (0.24,1.56) (-1.89,0.57) (1.15,-1.55) (-1.13,-1.56) (-0.32,2.06)
      (0.06,-1.62) (-2.24,0.96) (2.17,-0.24) (-2.09,1.01) (1.25,-2.05) (2.07,0.19)
      (-0.31,-1.51) (-1.78,1.32) (-2.07,0.01) (0.74,2.21) (1.84,-0.87) (0.72,2.20)
      (-2.31,-0.94) (-1.52,-1.06) (1.05,-1.30) (-1.85,0.39) (1.08,-1.98) (-1.92,-0.29)
      (2.01,-0.57) (-0.34,-1.59) (1.90,-1.17) (1.66,-0.63) (0.55,-1.60) (1.70,1.64)
      (1.46,1.45) (-1.48,-1.32) (-0.32,2.38) (-1.19,1.05) (1.01,1.95) (0.14,-1.82)
      (1.02,-1.33) (1.46,1.61) (-1.85,-0.20) (-1.45,1.12) (0.61,-2.37) (-1.66,0.37)
      (-0.24,-2.00) (-1.39,-0.61) (1.65,-0.23) (-2.09,1.15) (2.28,-0.18) (-1.77,1.04)
      (0.71,1.57) (0.41,-2.02) (-0.21,1.50) (-2.03,-0.89) (-1.35,-1.76) (0.67,2.04)
      (2.06,0.45) (1.57,-0.13) (0.72,-1.59) (1.48,1.45) (1.09,-1.55) (-0.13,2.49) (0.04,2.42)
      (0.23,-1.64) (0.10,-2.09) (1.25,-1.81) (1.07,1.24) (-0.03,-1.81) (-2.18,0.42)
      (-1.10,2.13) (-0.18,-1.83) (0.98,-1.44) (-1.02,2.09) (1.17,1.72) (1.97,-0.10)
      (1.53,-0.85) (-0.39,1.52) (0.60,-1.40) (1.76,1.11) (1.44,-0.88) (0.38,-2.45) (0.53,1.52)
      (-0.55,1.87) (-1.56,-1.07) (-1.07,1.36) (-0.19,-2.24) (-1.79,-1.18) (0.59,2.15)
      (-1.22,-1.01) (1.85,0.16) (1.54,1.30) (1.02,1.63) (-0.93,1.23) (1.69,0.13) (2.21,-1.04)
      (0.11,1.66) (-0.30,2.33) (-1.65,1.64) (1.76,-0.69) (-1.17,1.58) (-2.11,0.97)
      (-0.65,2.08) (2.31,-0.35) (-1.49,1.69) (1.90,1.10) (-1.32,-2.02) (-0.48,-2.27)
      (-1.17,1.21) (0.54,2.18) (-1.34,0.86) (-1.79,0.72) (1.90,-0.06)};
    \addplot[black,very thick] coordinates {(1.218,0.000) (1.216,0.064) (1.213,0.128) (1.211,0.192) (1.208,0.257) (1.206,0.323)
      (1.203,0.391) (1.199,0.460) (1.177,0.524) (1.140,0.581) (1.104,0.638) (1.068,0.694)
      (1.032,0.750) (0.996,0.807) (0.960,0.864) (0.922,0.922) (0.873,0.970) (0.810,1.001)
      (0.749,1.031) (0.689,1.061) (0.630,1.090) (0.570,1.120) (0.510,1.145) (0.446,1.162)
      (0.382,1.176) (0.319,1.191) (0.256,1.205) (0.193,1.220) (0.130,1.234) (0.065,1.243)
      (0.000,1.245) (-0.065,1.247) (-0.131,1.249) (-0.198,1.251) (-0.266,1.253) (-0.331,1.235)
      (-0.394,1.214) (-0.458,1.192) (-0.521,1.171) (-0.586,1.149) (-0.650,1.126)
      (-0.698,1.075) (-0.745,1.026) (-0.791,0.977) (-0.836,0.928) (-0.880,0.880)
      (-0.925,0.833) (-0.969,0.785) (-1.014,0.737) (-1.055,0.685) (-1.094,0.632)
      (-1.124,0.573) (-1.144,0.509) (-1.163,0.447) (-1.183,0.384) (-1.198,0.321)
      (-1.208,0.257) (-1.217,0.193) (-1.221,0.128) (-1.223,0.064) (-1.225,0.000)
      (-1.227,-0.064) (-1.229,-0.129) (-1.231,-0.195) (-1.233,-0.262) (-1.236,-0.331)
      (-1.236,-0.402) (-1.208,-0.464) (-1.172,-0.522) (-1.137,-0.579) (-1.101,-0.636)
      (-1.067,-0.693) (-1.031,-0.749) (-0.996,-0.807) (-0.954,-0.859) (-0.909,-0.909)
      (-0.861,-0.957) (-0.812,-1.003) (-0.763,-1.050) (-0.706,-1.087) (-0.641,-1.110)
      (-0.577,-1.133) (-0.514,-1.155) (-0.452,-1.177) (-0.389,-1.196) (-0.320,-1.194)
      (-0.253,-1.193) (-0.189,-1.191) (-0.125,-1.189) (-0.062,-1.188) (-0.000,-1.186)
      (0.062,-1.185) (0.124,-1.183) (0.187,-1.180) (0.250,-1.178) (0.315,-1.176)
      (0.381,-1.173) (0.446,-1.163) (0.511,-1.147) (0.574,-1.127) (0.637,-1.103)
      (0.701,-1.080) (0.767,-1.055) (0.834,-1.030) (0.905,-1.005) (0.966,-0.966)
      (1.021,-0.919) (1.071,-0.867) (1.117,-0.812) (1.164,-0.756) (1.188,-0.686)
      (1.194,-0.609) (1.200,-0.534) (1.206,-0.463) (1.211,-0.393) (1.216,-0.326)
      (1.221,-0.260) (1.225,-0.194) (1.223,-0.129) (1.220,-0.064) (1.218,0.000)};
  \nextgroupplot[mlplot,axis lines=left,width=7cm,xmin=0,xmax=500,ymin=0.005,ymax=1,
    ymode=log,xlabel={step},ylabel={training loss},
    xlabel style={at={(ticklabel cs:0.5)},anchor=north},
    ylabel style={at={(ticklabel cs:0.5)},anchor=south,rotate=90},
    title={training loss},title style={font=\small}]
    \addplot[PlotBlue,thick] coordinates {(1,0.8149) (2,0.6315) (3,0.5773) (4,0.5503) (5,0.5288) (6,0.5097) (7,0.4920) (8,0.4754)
      (9,0.4595) (10,0.4445) (11,0.4302) (12,0.4166) (13,0.4037) (14,0.3914) (15,0.3796)
      (16,0.3684) (17,0.3577) (18,0.3475) (19,0.3377) (20,0.3283) (25,0.2872) (30,0.2535)
      (35,0.2255) (40,0.2019) (45,0.1819) (50,0.1647) (55,0.1497) (60,0.1367) (65,0.1252)
      (70,0.1149) (75,0.1057) (80,0.0975) (85,0.0902) (90,0.0838) (95,0.0780) (100,0.0730)
      (120,0.0578) (140,0.0477) (160,0.0405) (180,0.0351) (200,0.0308) (220,0.0275)
      (240,0.0247) (260,0.0224) (280,0.0205) (300,0.0188) (320,0.0174) (340,0.0161)
      (360,0.0150) (380,0.0141) (400,0.0132) (420,0.0124) (440,0.0117) (460,0.0111)
      (480,0.0105) (500,0.0100)};
\end{groupplot}
\end{tikzpicture}
\caption{Left: the training data (blue: class $1$, red: class $0$) and the decision
  boundary learned by the network, where its output probability is $\frac12$. As a
  ReLU network computes a piecewise linear function, the boundary is a polygon.
  Right: the training loss, on a logarithmic scale.}
\label{fig:training-result}
\end{figure}

\begin{keytakeaways}[Where each idea came from]
\renewcommand{\arraystretch}{1.25}
\begin{tabular}{@{}>{\raggedright\arraybackslash}p{0.47\linewidth}
                >{\raggedright\arraybackslash}p{0.45\linewidth}@{}}
\textbf{Ingredient} & \textbf{Mathematics}\\
\hline
layers $\mW\va + \vb$, shapes, outer products & linear algebra (\cref{ch:linalg1,ch:linalg2})\\
activations, $e^x$, $\ln$, sigmoid & functions (\cref{ch:functions})\\
backpropagation, autodiff & chain rule, gradients (\cref{ch:calculus,ch:multivariable})\\
gradient descent, SGD, Adam & optimisation (\cref{ch:optimization})\\
initialisation, normalisation, attention scaling & variance (\cref{ch:probability})\\
softmax, cross-entropy, KL divergence & probability, likelihood (\cref{ch:probability,ch:statistics}, this chapter)\\
validation, regularisation, error bars & statistics (\cref{ch:statistics})
\end{tabular}
\end{keytakeaways}

\sectionexercises

\begin{exercise}
How many parameters does the network in the program have?
\end{exercise}
\begin{solution}
$\mW_1$ has $16\times2 = 32$ entries, $\vb_1$ has $16$, $\vw_2$ has $16$, and $b_2$ is one
number: $65$ parameters.
\end{solution}

\begin{exercise}
In the program, why is \texttt{d2} divided by $n$?
\end{exercise}
\begin{solution}
The loss is the \emph{mean} of the $n$ per-example losses, so the derivative of the
loss with respect to the logit of example $i$ is $\tfrac1n(q_i - y_i)$.
\end{solution}

\begin{exercise}
What would happen in the program if all the weights and biases were initialised
to $0$?
\end{exercise}
\begin{solution}
All hidden pre-activations would be $0$. The code uses $\ReLU'(0) = 0$, and
together with the zero output weights this makes every gradient of $\mW_1$,
$\vb_1$ and $\vw_2$ vanish: only $b_2$ would learn, and the network would
predict the same probability everywhere. More generally, identical initial
weights give identical neurons, which receive identical gradients and stay
identical; random initialisation breaks this symmetry.
\end{solution}

\begin{exercise}
What would you expect to happen if the learning rate were made much larger, say
$100$?
\end{exercise}
\begin{solution}
Steps would overshoot and the loss would oscillate or diverge instead of
decreasing, as in \cref{sec:learning-rate}; with ReLU networks, large steps can
also push many neurons into the inactive region permanently.
\end{solution}

%----------------------------------------------------------------------
\section{Chapter summary}\label{sec:ch11-summary}
%----------------------------------------------------------------------

\begin{keytakeaways}
\begin{itemize}
  \item A feed-forward network computes $\vz^{(l)} = \mW^{(l)}\va^{(l-1)} + \vb^{(l)}$,
    $\va^{(l)} = \varphi(\vz^{(l)})$; without non-linear activations it collapses to
    one linear map.
  \item $\softmax(\vz)_k = e^{z_k}/\sum_je^{z_j}$ turns logits into probabilities; it is
    shift invariant, so subtract $\max_jz_j$ before exponentiating.
  \item Entropy: $H(\vp) = -\sum_kp_k\log p_k$. Cross-entropy: $H(\vp, \vq) =
    -\sum_kp_k\log q_k$. KL divergence: $D_{\KL}(\vp\,\Vert\,\vq) = H(\vp, \vq) -
    H(\vp) \ge 0$. The cross-entropy loss is the negative log-likelihood.
  \item Matrix calculus: $\pderiv{L}{\mW} = \vdelta\vx^\top$, $\pderiv{L}{\vx} = \mW^\top\vdelta$
    for $\vz = \mW\vx + \vb$; activations multiply by $\varphi'(\vz)$ componentwise.
    Gradients have the shapes of the variables.
  \item Softmax with cross-entropy has gradient $\vq - \vy$: prediction minus truth.
  \item Backpropagation computes all gradients with one backward pass, using
    the recursion
    $\vdelta^{(l)} = \bigl(\mW^{(l+1)\top}\vdelta^{(l+1)}\bigr) \odot
    \varphi'\bigl(\vz^{(l)}\bigr)$.
  \item Products over many layers make gradients vanish or explode; ReLU, good
    initialisation, normalisation, residual connections and clipping help.
  \item Automatic differentiation records the computation and applies the chain
    rule in reverse.
  \item Attention: $\softmax(\mQ\mK^\top/\sqrt d)\mV$, a weighted average of values.
\end{itemize}
\end{keytakeaways}

You have now seen the mathematical core of modern machine learning, from counting
and fractions to backpropagation and attention. The last chapter suggests where
to go from here.
